\section{Evidence Revision and World Maintenance}
\label{sec:revision}

\begin{center}
\figfile{figures/figure4_revision.pdf}{\columnwidth}
\captionof{figure}{Evidence arrival and withdrawal differ from physical
building change. E2 and fresh F2 use the same effective sources; retained
knowledge and relation readouts are compared.}
\label{fig:revision}
\end{center}

The revision subset uses two buildings, eight configurations, and four
stages, totaling 64 stages. E0 receives plans; E1 adds a section
supporting an exact height; E2 withdraws it in the same context.
F2 constructs afresh from E2's effective evidence. The three sequential
stages are not independent repetitions.

\begin{center}
\small
\begin{tabular}{lrrr}
\toprule
Stage & Valid & Knowledge & Relations \\
\midrule
E0 & 15/16 & 15/15 & 68/80 \\
E1 & 10/16 & 10/10 & 49/80 \\
E2 & 13/16 & 13/13 & 58/80 \\
F2 & 14/16 & 14/14 & 65/80 \\
\bottomrule
\end{tabular}

\captionof{table}{Revision accounting. Valid: contract-valid outputs;
Knowledge: correct explicit status among valid outputs; Relations:
retained correct readouts over all allocated opportunities.}
\label{tab:revision}
\end{center}

All 52 valid stages have correct explicit height-knowledge states,
active sources, and support edges. All ten valid E1 worlds also
implement the supplied height. Withdrawal may return a measured
fact to unknown while leaving hypothetical completion coordinates:
those are different claims.

E2 and F2 have 11 jointly valid building--configuration pairs.
Their knowledge states agree 11/11, but retained readouts agree
51/55. E2 alone is wrong in three cases and F2 alone in one.
Explicit source-state accuracy and spatial retention therefore need
separate measurements. These differences do not by themselves
identify stale neural memory as their cause.

\section{Discussion and Limitations}
\label{sec:discussion}
\paragraph{One question, several observable capabilities.}
High R accuracy, a valid program, and a coherent fragment are different
achievements. Correct relations can improve construction without
ensuring valid interfaces or harmless additional material.
Qwen's N2 success also shows that the difference is not uniformly
large. The benchmark measures when these outputs align and diverge.

\paragraph{Scope of the evidence.}
The dataset is small and target-assisted. CS1 adds unit information;
N2 reuses v1 buildings; revision is not repeated under N2/CS1.
Public pretraining exposure and layout-family dependence are unresolved.
References describe selected drawing-supported relations, not complete
site measurements. Different provider tools and reasoning budgets
cannot isolate neural weights.

\paragraph{Scope of geometric checks.}
The compiler supports finite box/stair fragments, not general BIM,
arbitrary meshes, curved stairs, structural safety, or full human
clearance. Undeclared obstacles remain outside its scope.
The slab equivalence issue is an instrument limitation; future
corrections should be versioned and applied uniformly.

\paragraph{Scope of revision.}
Removing a source ID does not prove that information has disappeared
from candidate domains, bindings, initialization, or hidden state.
Identity and scope corrections need further task variants.
Future comparisons should distinguish them and include ordinary
reconstruction and source-tracking controls where applicable.

\section{Conclusion}
\tw{} evaluates whether answers, executable geometry, grounding,
and evidence updates remain compatible with one observed world.
The preserved experiments show condition-dependent task differences,
both benefits and regressions from correct relations, material
conflicts outside fixed queries, and a separation between explicit
knowledge updates and retained spatial conclusions.
Joint measurement gives a more specific account of spatial
understanding than answer accuracy or geometric validity alone.
